% backmatter/appendixB.tex —— 附录 B 可复现的自进化实验平台
% 本附录描述的是本书自己的实验台（`experiments/`），不是任何外部方法的介绍。
% 数字身份纪律：平台参数（预算上限、任务集规模、划分）是仓内配置，一律 \calcnum{}；
% 运行数与轨迹数是目录实测，一律 \measured{}；本附录不参与任何命题判决。
\chapter{可复现的自进化实验平台}
\label{app:lab}

本书的自有试验结论需要可追溯的运行，外部文献报告另列来源。实验台为此保存配置、轨迹、持久物和资源账。
本附录先说明历史运行的模块与记录口径，再介绍独立的离线组件。设计要求、已实现的接口约束与实际运行记录应分别核对；后加的检查不能追记成旧运行当时已经具备的保护。

\section{它要解决的三个问题}

\begin{itemize}
  \item \textbf{持久物必须真的持久。}一个 episode 结束后，除了被显式写入持久组件的东西，
        临时工作状态应按协议复位，轨迹可以归档。它对应第1章的状态边界要求：
        默认执行路径为每个 episode 新建临时目录并重新 materialize，正常返回前尝试清理；调用者传入沙箱或异常退出时，还须核对清理是否完成。临时目录重建也不等于进程与外部状态都已复位。
  \item \textbf{验收不能由做题者给出。}判定来自沙箱里的验收命令，不是模型自评。
        第 4 章的全部层级都建立在这条分界上，
        平台通过写工具保护指定文件；这不是完备的进程或操作系统隔离。
  \item \textbf{资源账要能分开记。}任务内花的钱与任务外（写记忆、渲染经验、评审）花的钱
        要分账，否则``同样的预算''这句话没有定义（§\ref{sec:appb-budget}）。
\end{itemize}

\section{一次运行的生命周期}
\label{sec:appb-lifecycle}

一次运行由 \texttt{run\_id} 标识，格式与落盘位置见附录 D。
平台内部把它分成五步：

\begin{enumerate}
  \item \textbf{载入配置。}任务集、划分、模型、温度、episode 配额 $E$、
        封存集重复次数 $R$、随机种子全部写进 \texttt{config.json}，
        并附上本仓库的 \texttt{git\_head}。
  \item \textbf{采集与写入（训练期）。}训练配额与写入规则按各实验的臂定义核对。以历史E01为例，四格各在训练集上花 $E$ 个 episode，完成后执行验收；只有格4在题后生成并写入持久经验，格1不写入，格2使用预置只读记忆，格3的适配笔记随当前任务结束丢弃。
  \item \textbf{判决期。}在预先划定的评价集上每题重复运行 $R$ 次，
        主指标 pass@1 由真实验收器给出。
  \item \textbf{聚合与判决。}按题配对求效应量与 95\% 自助区间，
        历史汇总另以已见方差估计区间半宽所需规模，不将其追溯称为预注册功效设计。
  \item \textbf{落盘。}每个 episode 一条 \texttt{trace.json}，
        运行级一份 \texttt{summary.json} 与 \texttt{config.json}。
\end{enumerate}

\section{模块职责}
\label{sec:appb-modules}

\begin{table}[htbp]
\centering
\caption{平台模块：每个部件只负责一件事，越界即违反平台纪律}
\label{tab:appb-modules}
\begin{tabular}{>{\raggedright\arraybackslash}p{3.9cm}>{\raggedright\arraybackslash}p{9.0cm}}
\toprule
模块 & 职责（以及它不负责什么） \\
\midrule
\texttt{backends.py} & 模型调用。只做真实调用与脚本回放两种后端，
        并记录瞬时失败重试次数。\textbf{不含任何价格}，美元一律留空 \\
\texttt{budget.py} & episode 级资源阈值与用量账。检查点发现超限时抛异常；
        \textbf{不控制进行中的请求，不内置价格表} \\
\texttt{tasks.py} / \texttt{taskset\_spec.py} & 任务定义与任务集规格。
        任务自带\texttt{protected}清单与族标签；规格独立保存正确与含缺陷实现，\textbf{生成的任务目录不带参考实现} \\
\texttt{build\_taskset.py} & 由规格生成 \texttt{experiments/tasks/v1/} 的目录与元数据 \\
\texttt{sandbox.py} & 临时目录、部分工具写保护、验收命令；旧路径检查有前缀漏洞，不是操作系统沙箱 \\
\texttt{static\_agent.py} & 本包保留实现含列目录、读写和调用测试；历史分派须按版本核对，不能只看提示清单。
        \textbf{它不学习}，持久更新由外层实验机制负责 \\
\texttt{trace.py} & episode 轨迹与运行标识。每条轨迹含判定来源与持久写入记录 \\
\texttt{eval.py} & 配对比较、自助区间、$Q$ 预算、封存集纪律、分族报告。
        \textbf{只出区间与原因，解释另行核对} \\
\texttt{e01\_arms.py} / \texttt{e02\_arms.py} / \texttt{e03\_arms.py} & 各实验的臂定义与经验渲染。
        \textbf{统计与判决须按运行身份回查原始档案} \\
\texttt{run\_e01.py},\newline\texttt{run\_e02.py},\newline\texttt{run\_e03.py} & 按臂安排回合、记账、汇总判定。公开源码收录 E01 入口；E02、E03 的运行器保存在完整维护档案。各次历史汇总均可由附录~\ref{app:repro} 的输入清单定位。 \\
\texttt{trigger\_signals.py} & 三个更新时机的触发信号计算（第 13 章） \\
\texttt{check\_backend.py} & 后端就绪探测；\textbf{不打印密钥} \\
\bottomrule
\end{tabular}
\end{table}

\section{文件工具的保护范围与实际缺口}
\label{sec:appb-sandbox}

旧平台为每个 episode 建立临时工作目录，验收子进程仍使用当前用户的权限。若要执行来源不受信任的候选程序，应在独立的隔离环境中运行，并限制它能读取的文件、使用的网络和服务凭证。任务声明的\texttt{protected}集合由指定写工具检查，直接修改匹配测试文件名的请求会抛出\texttt{SandboxViolation}。执行器把拒绝写入\texttt{tool\_result}，所以不能仅凭步骤类型将其当成成功动作。这个集合没有改变操作系统文件权限，也不控制任意子进程的标准库读写。

\begin{sloppypar}
\textbf{工具规则与进程权限是两层不同的边界。}2026年9月30日的临时文件检查记录了旧工具对目录归属的误判，也记录了普通子进程能够修改受保护标记。历史报告保留了这些观察及对应源码的校验值。当前文件工具按目录层级检查访问，候选程序还需要独立的受限执行环境；本版在缺少该环境时停止调度。
\end{sloppypar}

新版E01文件工具采用公开文件清单和可写子集，按解析后的父子路径关系检查归属，并拒绝歧义路径、设备名、硬链接和已识别的链接路径。配套离线检查覆盖相邻目录、未列明文件、公开测试写入和输入分区。这些修复位于独立文件，不追溯修改已冻结实验依赖。它们仍不是对恶意并发文件替换的保证，也没有把普通Python执行变成操作系统隔离；E01的真实候选运行因此仍有明确阻断项。

历史数据的处理也要有边界。发现实现存在可绕过路径，不等于证明每条旧轨迹都实际利用了它；反过来，旧工具日志中没有越界记录，也不能证明绕过工具的写入从未发生。当前保留历史成绩和已观察行为，将完整性证据不足列为限制，不删除失败或重算出一套看起来更干净的旧数据。若新版本改变执行权限，应用新运行编号和清楚的版本说明。

验收命令默认用标准库 \texttt{unittest}。
这是当前任务集采用的执行入口；历史环境中曾记录 pytest 安装受限，不能据此断言今天仍不可安装。
更换验收框架需要核对发现规则、退出码与隔离方式，验证行为一致后再比较结果。

\section{预算、副账与成本纪律}
\label{sec:appb-budget}

\begin{table}[htbp]
\centering
\caption{episode 级预算的默认值（\texttt{budget.Budget} 的字段）}
\label{tab:appb-budget}
\begin{tabular}{>{\raggedright\arraybackslash}p{5.2cm}c>{\raggedright\arraybackslash}p{5.4cm}}
\toprule
字段 & 默认 & 当前实现的限制 \\
\midrule
\texttt{max\_calls} & \calcnum{8} & 成功返回的模型调用；网络重试另计 \\
\texttt{max\_steps} & \calcnum{12} & 执行循环次数；格式错误也消耗一步 \\
\texttt{max\_prompt\_tokens} & \calcnum{200000} & 输入 token；收到回复后才可观察 \\
\texttt{max\_completion\_tokens} & \calcnum{20000} & 输出 token；不能在流式生成中强制中止 \\
\texttt{max\_wall\_seconds} & \calcnum{600} & 循环入口检查墙钟；不是进行中请求的硬截止 \\
\bottomrule
\end{tabular}
\end{table}

预算检查发生在执行循环入口。下一次检查发现超限时，\texttt{BudgetExceeded} 抛出，循环停止，随后仍进行最终验收并记账。进行中的调用可能使 token 或墙钟越过阈值；若该回复直接给出\texttt{FINISH}，程序不会再回到循环入口，因此也可能没有预算异常标记。没有异常不等于没有超额，用量应按实际记录核对。

一个离线教学夹具把输入 token 阈值设为\calcnum{100}，再让预写回复报告\calcnum{110}的用量。回复直接结束时，原循环保存\calcnum{110}而不产生预算异常；若同样回复要求继续使用工具，下一轮入口才拒绝继续。夹具使用构造用量、固定时钟和固定验收对象，没有模型请求或候选程序执行。它检查控制流，不提供模型成本或能力证据。

旧 E01 运行器的\texttt{--trial}曾以较少的题目和重复次数调用真实后端。本版同名选项仅打印离线计划。要核对当时的用量，可运行\path{experiments/lab/cost_report.py}读取归档汇总。旧轨迹记录成功回复数与token用量；其中episode墙钟在最终验收前取值，后续验收、落盘与清理的耗时应另计。

\begin{sloppypar}
任务外的模型调用（写记忆、渲染经验条目、学习式评审）计入\textbf{副账}
（\texttt{side\_calls} / \texttt{side\_prompt\_tokens} / \texttt{side\_completion\_tokens}），
历史E01/E02不占episode配额，但包含在各臂总调用与token里；E03来源、渲染与测试账覆盖不同，不能套用这一合计约定。新E11/E12按各自协议实施共同预算。
分账的理由来自公共协议第 6 条：重试、候选筛选、模型评审与失败运行都要进资源账。
不分开记的话，``两臂预算相同''这句话无法验证。
\end{sloppypar}

\begin{sloppypar}
\textbf{美元一律不填。}当前运行器未接入可靠计价表，\texttt{Usage.usd} 恒为 \texttt{None} 并置
\texttt{unpriced=True}。本书所有成本比较因此只报实物量（调用数、token、墙钟），
折美元要等 P-1 的 \texttt{calculations/pricing.yaml} 拿到有出处的定价表。
正文里出现的美元数字全部是\emph{引用}自他人论文，不是本书实测。
\end{sloppypar}

\section{评测口径}
\label{sec:appb-eval}

\begin{itemize}
  \item \textbf{配对比较}：当前\texttt{pair\_up}先检查两臂的\texttt{task\_id}各自唯一、运行判定非空且为二值，并核对同一任务的族与划分。单侧缺项双向列出，配对取共同任务。历史E01使用的旧helper只有单向缺项提醒，调用者还须检查重复标识与空运行。完整配对分析须明确筛选范围与偏倚风险，缺失如何计分由预定估计目标决定。
  \item \textbf{不确定性}：95\% 百分位自助区间，重采样单位可选``题''或``任务族''。
        两种单位对应不同的推广范围，区间宽度的次序也未必固定；不能把相关题目或重复调用直接当作更多独立任务族。
  \item \textbf{只报效应量与区间}，不报裸 p 值。
  \item \textbf{$Q$ 预算}：辅助检查器可以约束查询预算，但调用者仍须正确接入；封存数据不得参与本轮候选产生与选择。终评查询与后续改动须留痕。
  \item \textbf{分族报告}：避免一个总分遮蔽局部退化。
  \item \textbf{精度与功效分开}：历史工具中的 $n\ge(1.96\,s/h)^{2}$ 近似控制区间半宽 $h$，不是包含检验功效的确认性样本量。正式规划见第\ref{ch03}章与附录A。
  \item \textbf{判定遵循对应协议}：规模、停止规则、比较对象和信息边界应共同满足要求。历史精度规划不冒充预注册门槛，也不能只凭达到某个整数就自动改判。
\end{itemize}

\section{任务集与划分}
\label{sec:appb-tasks}

\begin{sloppypar}
任务集 \texttt{v1} 为 \calcnum{8} 个任务族 × \calcnum{8} 道题，共 \calcnum{64} 道；
四个用途划分（train / select-valid / regression / sealed）各 \calcnum{16} 道，
\emph{每个划分里 8 个族各 2 道}。
八个族是：边界与差一、可变默认参数、字符串索引、字典缺键、
类状态、错误契约、递归基、排序键——每一族对应一类具体的、可命名的编程错误。
\end{sloppypar}

\begin{principlebox}
\textbf{迁移范围按实际来源界定。}
train 与 sealed 目录覆盖同一组族，不表示某次经验构建实际用到了全部族。
评价题属于来源已见族还是未见族，应逐次运行核对；若两类同时存在，整体读数便混合两种范围，
不能只凭目录名解释为纯同族或纯跨族迁移。
E03另设\emph{族留出}拆分，整两个族不参与经验构建，测试取这两族的题。
该比较范围限于所留族，也未交叉操纵形式与容量。历史做法见 \texttt{experiments/protocols/E03.md} §3；
拆分A基线冲突的采用范围见表~\ref{tab:ch05-e03a}。
\end{principlebox}

E01 类文本经验检索刻意\emph{不用}任务族标签：族是任务集的元数据，真实 Agent 拿不到，
用它检索等于把标签泄漏进记忆索引。这一检索器用字符二元组重叠，取前 \calcnum{2} 条，
重叠全为 0 时退化为取最近的两条。E05 按任务族组织技能序列，是已知族条件下的受限原型，不能把前述无族标签声明套到 E05。

\clearpage
\section{已完成的运行}
\label{sec:appb-runs}

\begin{table}[!ht]
\centering
\caption{已归档的模型运行与早期演示（离线组件另见后文）}
\label{tab:appb-runs}
\begin{tabular}{>{\raggedright\arraybackslash}p{5.0cm}>{\raggedright\arraybackslash}p{1.6cm}>{\raggedright\arraybackslash}p{6.4cm}}
\toprule
运行 & 实验 & 说明 \\
\midrule
\texttt{E01-20260923T152230Z}\newline\texttt{-90acda} & E01 & 预算分配四格；\measured{352} 个 episode，四格结果出齐，仍须保留历史设计局限 \\
\texttt{E02-20260924T153511Z}\newline\texttt{-421eac} & E02 & 反馈方式六臂试跑；每臂 \measured{88} 个 episode \\
\texttt{E03-20260925T152453Z}\newline\texttt{-254845} & E03 /\newline E03c & 两个拆分共\measured{10}臂；\measured{800}个测试episode，最后一臂跨日期补齐 \\
\texttt{E05v2-20260929T143245Z}\newline\texttt{-d7a8db} & E05 v2 & 单族\measured{6}序列、\measured{9}主episode；工程流程试跑 \\
\texttt{CH02DEMO-20260923T}\newline\texttt{153152Z-56208b} & 演示 & 第 2 章章首的状态残留演示（不是实验） \\
\bottomrule
\end{tabular}

\vspace{2mm}
\begin{footnotesize}
运行标识格式见附录 D。每次运行的原始记录在 \texttt{experiments/runs/<run\_id>/}，
已有渲染副本位于\texttt{experiments/results/}，但该目录不完整；本版离线入口按\nolinkurl{analysis/inputs.json}读取。汇总与轨迹冲突时先溯源有效记录集，不能仅凭目录位置裁定。
重新执行与续接旧运行是不同操作。新版E05/E11/E12拒绝覆写既有运行目录；不能将书稿构建的清理开关当成实验续跑规则。
\end{footnotesize}
\end{table}

\section{把一条汇总结果追溯到原始轨迹}
\label{sec:appb-reconcile}

运行目录存在，并不意味着里面的数字已经可以用于正文。
本书建议按配置、序列、轨迹、汇总的顺序核对，先确认这些对象是否描述同一次执行，再讨论结果是否支持命题。
这个顺序能够发现一种常见错误：汇总文件来自后一次重跑，轨迹目录却混有前一次失败留下的文件。
文件都真实存在，拼在一起仍可能产生不存在的一次实验。

配置检查先核对任务标识、划分、模型身份、温度、预算和实现版本。
只保存 Git 提交号还不够：运行时可能有未提交修改，任务文件或协议也可能没有被版本控制覆盖。
E05 v2 开始另外保存实际源文件、协议和任务的校验信息；它们帮助确认内容是否相同，不承诺外部模型服务在未来仍然相同。
固定模型别名与固定模型权重不是一回事，记录一个随机种子也不等于后端实际接收了它。

序列检查回答哪些任务共享持久状态。
同一族、同一复用格点、同一次重复中的构建与后续任务构成一条序列；另一格点重新建立独立状态。
因此不能为了省事把较小格点的轨迹复制到较大格点，再称为一次新的独立运行。
如果另一个实验选择共享前缀，那是可以单独定义的设计，但必须同时处理结果相关性与预算复用。

轨迹检查把模型调用、工具动作和验收分开。
模型回复包含一段看起来正确的代码，不代表代码已经写入；写工具返回拒绝，也不能被当作成功写入重放。
只有工具成功执行后的内容才是候选工件的一部分，最终结果则来自实际验收。
来源代码提取如果忽略这层区别，就可能把模型想写的程序和真正被测试的程序混在一起。

成本合账随后逐项进行。
先在每个 episode 内检查调用与 token，再按构建、复用和额外修复归入各自阶段，最后按序列和实验臂相加。
一个已经合计了全部复用任务的字段，不能在绘制成本曲线时再次乘以复用次数。
这类错误不会产生明显的程序异常，却会系统性改变交点的位置；因此核算应从最细的可得账目向上进行。

同一份账里还可能存在不同时间范围。
进程启动到结束的墙钟包含准备、模型等待与检查，episode 计时可能只覆盖其中一段，后端请求计时又更窄。
这些量应以各自名称报告。将 episode 秒数相加后称为整次实验墙钟，或将续跑进程的重试计数称为全程重试数，都会改变读者对资源的理解。
没有保存完整时间信息时，应报告范围缺口，不能用已有局部数填补全程。

最后检查汇总是否保留失败与缺失。
任务失败通常仍有成本；网络中断可能有成本却没有可观测 token；某个序列从未开始则没有任务结果。
三者不能都填成零。只保留完整配对的分析也要报告排除了哪些对象，以及缺失是否可能与任务难度或实验臂有关。
配对筛选是一项分析决定，不是把不完整数据从研究记录中删除的理由。

\section{E05 小规模试跑提供了什么平台证据}
\label{sec:appb-e05-smoke}

2026年9月29日，E05 v2 在一个任务族、两个训练题上完成了一次授权工程试跑。
三个臂、两个复用格点形成 \measured{6} 个序列，共 \measured{9} 个主 episode、\measured{48} 次模型调用。
逐轨迹核算与运行汇总相符；没有额外修复 episode，当前进程后端重试计数为零。
这里登记的是平台执行范围，能力解释见第\ref{ch08}章的限制。

静态技能组完成来源准入，并在一个后续任务中读取参考。
动态组两个来源任务则在格式解析处停止：一个回复使用另一种工具标记，一个回复在单行 JSON 约定内包含实际换行。
这两条来源均未通过验收，准入层因此没有发布技能。
可确认的是拒绝链路与成本记录工作了；不能确认动态修订在真实模型下有效，因为此次运行没有走到那条分支。

这种结果说明，工程试跑的完成条件应包括暴露错误，而不只是让一个汇总文件成功落盘。
若把两次提前失败造成的少调用写成效率优势，就会将装置兼容性问题转写为方法优势。
若删掉失败任务再计算成功率，又会改变原先授权并执行的样本范围。
本书保留全部轨迹，并把它们作为后续离线修复的输入。

新增的 v3 执行器保留原始回复，同时单独记录解析后的动作。
对于已观察到的格式，可以在不更改候选代码内容的前提下解析；对于多动作、重复字段或未知工具，则拒绝执行并请求格式纠正。
纠正使用同一 episode 剩余调用预算，不能每次报错重新获得额度。
原始回复、拒绝原因与实际执行的动作各自留痕，才可以检查兼容处理有没有悄悄替模型补做任务。

离线检查包括两条真实失败回复的解析回放、格式错误连续出现时的调用上限、受保护测试的写入拒绝，以及动态修订的完整状态转换。
在脚本后端驱动的例子中，来源成功、后续首轮失败、修复成功、再一次失败被拒绝额外修复的过程可以被逐项核对。
该过程说明软件按给定回复运行，不是模型通过了这些任务。
v3 尚无真实运行结果，不能把离线过程的通过率列进模型实验表。

\section{中断之后怎样决定能否续跑}
\label{sec:appb-resume}

续跑首先要回答的是已有记录是否完整，而不是目录是否存在。
一个序列有完整结束标记、所有轨迹都可读、来源哈希一致，才有条件作为已完成单位复用。
只有开始记录却没有结束标记的序列，可能已经发送请求，也可能已经执行工具；没有足够信息时，不应直接从同一路径重写。
否则原有消耗与外部效果会被隐藏在一次看似干净的新尝试之后。

E05 新版对未完成序列采取保守处理：保留当前工件，停止自动覆盖，要求先核对已发生的动作和费用。
这不是完整的断点恢复系统。它没有承诺能够自动恢复任意工具中断，只保证不把已知不完整的序列悄悄当作全新序列重跑。
未来若加入细粒度恢复，需要进一步记录动作是否已执行、执行是否可重复、持久状态是否已提交。

已经完成的序列也可能因条件变化而不适合继续拼接。
后端模型变更、提示修改、预算定义改变、任务划分修订，都可能改变被估计的问题。
这时应使用新编号，明确新旧版本的关系；保留旧记录能够支持对照审计，却不授权把它们自动合并为同一批数据。
即使只修复了一个看似小的解析问题，也可能改变任务成功率与调用分布，因此应说明修改发生在看到哪些结果之后。


\section{新版时机与恢复实验的离线入口}
\label{sec:appb-v2entry}

E11 v2 与 E12 v2 使用独立文件和新运行目录，保留旧实现与旧记录。两者的默认入口计算并打印预检清单，不构造真实模型后端；需要保存时，另用\texttt{--save-preflight}指定尚不存在的输出路径。清单包括当前协议、实现依赖、选中任务及后端身份，目的是检查实际运行内容是否与此前核对的版本一致。保存清单不会发生模型请求，也不代表已经得到付费运行授权。

完整维护库包含\path{experiments/tasks/v1/}任务目录。若已有任务，先核对其版本，再从仓库根目录运行以下只读预检。公开预览初版只提供生成规格，首次预检前须先按下段生成任务。
\begin{verbatim}
python -m experiments.lab.run_e11_v2
python -m experiments.lab.run_e12_v2
\end{verbatim}

若某个源码包确实不含任务目录，应在独立副本中按该版本规格生成，再核对输入身份：
\begin{verbatim}
python -m experiments.lab.build_taskset --no-lab-run
\end{verbatim}
这个命令仍会计算规格中的参考与缺陷样例，并覆写目标目录的题面、代码、测试和元数据；该开关只跳过后续实验台实跑，不表示只读。不要用它刷新历史任务输入。E11/E12默认读取仓库内的固定任务路径，生成到另一个\texttt{--out}目录也不会自动改变预检输入。

真实执行入口还要求显式执行开关，以及与当前预检逐项一致的批准清单路径。这个文件检查只能发现版本变化，不能判断用户是否真的同意付费；操作者仍需核对已有授权覆盖的任务、调用规模和数据范围。把文件命名为“已批准”不会自动赋予授权，旧版的一次试跑额度也不会因复制清单而变成新版额度。

E11 v2 将每题做题和经验构建纳入同一序列调用账，时机策略仅能在共同配额内决定是否尝试构建。它保存当前题验收、监测信号、构建请求、候选文本和最终发布状态。末题之后不再构建，失败或格式拒绝消耗已申请的构建机会。脚本后端的完整序列检查可以确认这些规则是否执行，却不能证明事件触发策略的收益优于其他策略。

E12 v2 从固定来源题构建共同初始记忆与候选池，评价题在任何来源运行之前固定。每个臂都从同一个初始版本克隆，不能继承上一臂的状态。四个更新臂均安排恢复与复评，使未触发业务回退线的臂也有恢复记录。共享准备、初始评价、各臂评价与恢复分别保存费用范围，再核对总和，不把同一份准备费复制成多个真实支出。

两者都在发送请求前写入原始消息与待返回条目，返回后补记原始回复、用量与状态。若进程在发送后中断，待返回记录说明请求可能已经派发；它不是零费用证明。新目录拒绝覆盖，也没有自动从中断位置续跑的承诺。恢复工作必须先区分请求是否发生、工具是否执行、结果是否保存，不能仅凭汇总文件缺失便从头再付费。

E11/E12的装置检查覆盖本地执行、预算中断、未知用量、候选拒绝、目录不覆盖与整段合账。检查索引和历史数量保存在版本化补充材料中；本附录按实际控制边界说明功能，真实模型兼容性、策略收益和恢复后行为仍未取得对应数据。

冻结清单也有覆盖边界。它可以核对本地文件内容，没有固定远程模型权重，也没有消除服务时段变化。当前实验的种子是记录标签，不能据此声称提供方生成过程完全可重现。做题工具可以读取测试文件，路径写入保护也不是操作系统级隔离；这些限制必须随结果一起传播，不能在复现实验时省略。

\section{平台不做什么}
\label{sec:appb-limits}

这一部分与功能同样重要——读者要能判断本书的哪些结论\emph{不可能}由这个平台得出。

\begin{itemize}
  \item \textbf{不训练任何参数。}平台里没有任何梯度更新，
        $\theta$ 从头到尾是黑盒 API。第 10、11 章涉及的 RL 与参数更新
        在这个平台上没有对应物。
  \item \textbf{不做真实工具生态。}本包保留的静态执行器含列目录、读文件、写文件、回合内验收与结束动作；历史运行是否使用某个动作，须核对当时实现与轨迹，不能仅按“旧版／新版”归类；
        没有专门的包管理、网络访问或多 Agent 通信工具；这不表示宿主上的候选进程已被禁止联网。第 9 章的多 Agent 结构
        在这里无法实测。
  \item \textbf{不做长程任务。}episode 上限是 \calcnum{8} 次调用、\calcnum{12} 步，
        更长的软件工程任务可能超出这一范围。本书的实测结论不能据此外推到长程。
  \item \textbf{不产生货币账单。}见 §\ref{sec:appb-budget}，美元一律留空。
  \item \textbf{统计输出与结论解释分开。}运行器提供区间与原因，结论还需结合问题、比较设计与适用范围。
\end{itemize}

\begin{sloppypar}
这五条合起来划出一条清楚的边界：
\textbf{本书的实测结论只覆盖``短程、单 Agent、以文本或轻量代码为持久物、
判定来自可执行验收器''这一象限}。
本书在这一范围提供自己的任务与轨迹，不能因此声称既有研究在该范围完全空白；
这些轨迹仍带有本附录列出的执行与信息边界，不能据此覆盖长期训练和真实工具生态。
\end{sloppypar}

\section{走查一条不会产生能力结论的离线流程}
\label{sec:appb-e01-offline}

E01新版目前有一套独立的离线组件，用于验证正式实验之前的控制与记录行为。它没有付费后端入口，也不运行模型生成代码。流程中的候选文本、公开分数和最终结果均为人为构造，不能把文件里的成功标记统计成Agent正确率。这里把它列为平台案例，是为了说明怎样检查测量装置，而不是补上P1的缺失结果。

\subsection{把同源起点保存为实际字节}

离线流程先创建一份手写的初始记忆$M_0$，再分别建立R、F、U三个臂的独立目录。R表示无记忆的重试，F表示只读冻结记忆，U表示持续更新。F与U拿到相同的不可变记忆对象和相同字节哈希；R的实际输入记忆为空，但运行清单仍记录准备物的身份，以便解释方案之间的信息差异。

每个臂执行\calcnum{2}次从头开始的重复，每次包含\calcnum{2}个合成任务机会，所以共有\calcnum{6}条序列、\calcnum{12}个任务机会。这些数量是离线程序的配置与事件计数，不是独立统计样本。第二次重复重新从$M_0$开始，不读取第一次U臂末尾的记忆；否则它就不是重复，而是延长了一条学习轨迹。

U臂完成首题后，才允许对来自该题的候选经验作发布决定。还没完成的当前题、尚未到达的未来题和更早但已错过发布窗口的题，都会被状态机拒绝。拒绝准入的经验不改变记忆哈希，发布成功才生成新记忆。末题后没有预定消费者，因此不再写入；若协议明确安排冻结后的迁移评价，才可以在相应配置中开放这次更新。

这一检查解决的是“什么时候、从哪一条记录写入”的工程问题。它不证明经验内容为真。当前\texttt{validated}是调用者传入的准入决定，未来真实实验必须定义谁作这个决定、看哪些材料以及验证消耗多少资源。不能把一个布尔字段命名为“验证通过”，就当成已建立语义正确性证明。

\subsection{一次失败怎样留下两本都能核对的账}

每个合成任务有\calcnum{3}次逻辑生成机会、最多\calcnum{4}次请求尝试、\calcnum{3}次公开测试机会。F和R把逻辑生成机会都用于求解候选；U在首题用其中\calcnum{2}次求解，另\calcnum{1}次形成更新。这样比较的是同调用上限下的资源分配策略，不是“额外免费写一次记忆”。U的末题则把\calcnum{3}次都用于求解。

脚本故意让U首题的更新请求先发生一次模拟传输失败，再返回预设文本。一次逻辑生成因此对应两次请求尝试。账本在调度前先写入请求事件，失败后记录结果未知的费用；成功返回的token计数只作为脚本给定值保存。它没有理由推断那次失败在真实服务里一定免费，更不能把失败从请求次数里删除。

完整流程因此形成\calcnum{36}个逻辑生成事件、\calcnum{38}个调度事件和\calcnum{34}个公开测试事件，两次模拟失败的用量保持未知。这组数可以从每题日志逐项相加：U少生成了两个求解候选，却多出了两个更新生成；两次更新各多一次模拟重试。没有任何一个事件向外部模型发出请求，实际API调用数为零。

逻辑调用、网络尝试和模型成功回复是三个计数器。把它们压成一个“调用数”字段会丢失故障信息：预算可能按逻辑生成配额分配，服务限额或重试风险却与调度次数有关；成功回复只覆盖可见的返回，不能说明其他请求没有到达服务端。结果报告必须同时给定义和合账关系，不能在不同图表里悄悄换用分母。

\subsection{提交必须早于最终评价的读取}

脚本候选中，第二个候选的公开得分最高，但预设最终结果是失败；其他候选中存在预设最终成功者。选择器只收到公开候选对象，仍然选第二个，并先保存提交文件。这个反例是故意构造的：如果系统最后报告它选择了隐藏结果成功的那个，就暴露了使用不可用评价信息的路径。

选择规则还有两个容易漏掉的边界。公开测试未运行的候选排在已运行候选之后，即使已运行候选得零分；相同公开分数按产生顺序取最早者，不按列表当前排序挑选。候选的公开测试总数必须一致，否则“通过项数更多”可能只是测了更多项。所有候选都不可提交时，选择结果应为空，而不是凭空生成一个默认成功答案。

这里得到的是接口和调用顺序证据：选择函数不接收最终结果字段，提交先落盘，更新使用已完成任务的公开材料。它不意味着任意恶意Python代码无法读取宿主上的评价文件。后者需要进程权限、环境与存储边界；用同一个解释器里的不同字典表示公开和隐藏，只能防止一部分意外混用。

\subsection{日志能重放，不代表请求能自动重发}

资源日志给每个事件一个身份。相同事件再次被读取，不再加一次费用；若身份相同但内容不同，则拒绝合并。这可以阻止恢复程序把已经记过的准备费用重复累计。读取完整日志后，恢复出的计数应与先前汇总一致；尾行截断、字段冲突或预算配置改变时，程序停止并保留原字节，交给后续核对。

更难的是发送后、返回前进程中断。日志只能知道曾为这次派发留出事件，不一定知道服务端是否已经执行和计费。当前实现遇到这种在途状态就阻止盲目重发，既不把费用清零，也不自动补一个失败结果。它没有声称外部请求“恰好执行一次”。只有服务端支持相应的请求身份与查询语义，或人工/独立流程取得了足够证据，才可能进一步恢复确定状态。

这种拒绝恢复可能牺牲便利，但保留了实验的解释边界。若直接删除在途事件再跑，新的成功结果会遮住第一次执行的不确定费用；若直接把在途算失败，又可能把系统观察故障当成模型失败。原始记录、派生分析和新的真实运行应分别编号，不能让一个“继续”按钮悄悄改变证据来源。

\subsection{共享准备物只生产一次，独立部署却各有取得成本}

离线账本支持两种明确分开的汇总。研究实际总账按准备账本身份去重，因此同一份准备物被多个臂引用时只计一次。若同一身份对应不同日志内容，即便数字总额相同也拒绝合并，以免相同金额掩盖不同执行事件。独立部署表则为每个需要$M_0$的方案附上取得准备物的成本，说明该方案离开研究共享环境后要承担什么。

本次$M_0$是手工写好的测试夹具，没有模型准备请求；这只意味着API事件为零，不意味着人的编写、检查和设计时间为零。当前没有测量这部分工时，就保留未测量说明。将来若$M_0$来自真实模型生成，应接入单独准备账本，再复用相同的共享身份；不能把手工夹具的零API记录复制成真实准备成本。

\subsection{平台通过了哪一层检查}

E01新版的装置检查围绕四条数据关系组织：起始记忆同源且不可变，候选与选择记录相连，在线和迁移身份互斥，调用与费用只能结算一次。对应反例覆盖越界文件、预算耗尽、中断重放与缺失判定；历史检查计数留在补充材料。数量本身不表示覆盖率或模型成功概率。

清单把准备、在线、冻结迁移分为互斥的题目身份，并要求声明的迁移族不与准备及在线族重合。初始记忆中的来源必须属于准备分区，分析行从同一清单展开，避免运行使用一套分区、分析再凭结果选另一套。内容摘要不包含题号和族名，因此只换题号不能使同一内容变成新题；与调用者提供的历史内容摘要相同的题也被拒绝。当前仅用构造材料检查这些条件，尚未提交一套新的确认性任务。

这套检查仍有明确盲区。语义相近的两道题、同一模板改写、未申报的历史暴露，不一定得到相同摘要；把一个族改名，也不会让它在统计上独立。因此清单哈希用于固定可审查的对象，不能替代任务设计、来源核查或语义去重。公开负载会排除声明为最终测试的文件，但如果作者误把答案写进公开提示，文件分区本身也无法替作者识别这项泄漏。

真实候选执行、公开与最终测试的技术隔离、新任务族与封存数据、语义准入规则、真实后端费用适配以及层级样本规划仍待完成或接通。这些缺口分别影响工程、信息边界和推断范围。即使全部离线检查持续通过，研究问题所需的模型能力证据也不会由此产生；只有符合新协议的真实观察，才能进入相应的能力证据链。


\section{把候选输出与验收权力分开}
\label{sec:appb-separate-oracle}

一个候选程序可以打印“全部通过”，也可以把正确答案混进报错文本；验收器不能把这些文字当成自己的判定。E01新版新增的离线组件先固定一个更窄的接口：每次调用只传一个JSON输入，候选返回含\texttt{result}的JSON对象，正确答案由外部比较器保管。这里定义的是JSON函数任务接口，尚未把旧的\texttt{unittest}题库转换成这种任务，更没有悄悄替换旧实验的分数。

候选包从任务的公开文件构造，替换只能发生在声明的可写文件中。最终测试文件不进入包；当前输入进入候选，其他测试输入、期望输出、验收代码与既往判定不进入这次调用。最终输入仍会被待测程序看到，这是函数求值所必需的；要阻止它把这些输入写回学习过程，还需要固定候选、关闭最终反馈回流，并在测试之间重建执行环境。

基础评测实例记录候选、套件和发布状态的摘要。一个测试只能发放一次，返回结果必须对应同一调用身份和同一候选；重复返回或替换候选被拒绝。调用尚未完成时报告缺失，不能让它从分母消失。这个基础对象只有内存状态；本节后面的持久日志进一步记录派发和恢复，但记录发布快照仍不等于在真实 Agent 中完成部署。

\subsection{先分清谁失败，再计算分数}

程序输出不符合约定，是候选失败；容器服务没有启动成功，是测量没有发生。这个区分直接影响分母。新版构造检查将启动失败、传输失败及清理状态未确认记为缺失；候选异常、已执行后的超时、内存或输出限额记为观测失败。状态必须由可信执行适配器记录，不能由候选自己宣称“我是服务故障”来免除失败。

外部比较器拒绝重复JSON键、非有限数字、额外的通过标记，以及合法对象后附带的其他输出。当前比较规则是规范JSON字节相等：\texttt{true}不等于\texttt{1}，\texttt{1}与\texttt{1.0}也不同。这是明确而狭窄的任务合同，适合要求确切结构与类型的题；若题目需要数值容差，必须预先定义容差和比较器，不能在看到错误答案后放宽。

公开测试的反馈只返回计划项数、已观察数、通过数、缺失数，以及由这些计数确定的“是否全部通过”状态。最终测试评测实例则拒绝生成供模型使用的反馈，原始标准输出和错误输出也不进入更新材料。候选仍可能猜中部分答案，测试本身也可能覆盖不足；这些是评价效度问题，不会因协议字段清晰而自然消失。

\subsection{配置限制必须落实为进程限制}

离线组件声明了输入、输出、源码大小以及墙钟、内存、进程数上限，但尚无运行时来执行后面三类限制。收到完整字节后检查长度，也不等于在读取过程中限制内存。Python官方文档说明，\texttt{communicate}会把读取数据缓存在内存；其超时处理与\texttt{run}的行为也须区分，不能把某个接口会终止直接子进程推广成整棵进程树必然被清理\cite{python-subprocess-docs}。

容器配置提供了实现隔离的一部分手段，例如网络设置、只读文件系统、资源限制及权限约束；这些选项的具体含义可查Docker官方文档\cite{docker-run-security-docs}。本书当前环境尚未验证一套可用配置，也没有把“写出参数清单”当成已经完成的保护。真实入口仍明确拒绝执行候选，没有退回普通宿主进程的选项。

下一层验收必须实际尝试读取虚构的宿主标记、写入只读候选、建立网络连接、超出资源上限并留下子进程，再检查每一项是否按预定规则被阻止和清理。探针必须使用新建虚构材料，不能拿真实凭据作测试。即使探针通过，也只说明已覆盖的配置与故障路径得到验证；未知漏洞、系统更新和未覆盖接口仍需要独立维护。当前\calcnum{29}项输入输出检查只运行构造字节和状态对象，没有启动候选进程，不能记为能力样本或系统隔离验收。

\subsection{先保存派发记录，再允许执行}

内存中的“已运行”标志无法跨越进程退出。假设一个测试已经发给执行器，控制程序却在收到结果前中断；重启后若只看结果文件是否存在，就可能再次发出同一测试。即便候选固定，多一次执行也会改变随机环境下的观察机会；若允许据此换候选，最终评价还会进入选择过程。没有结果文件，只能说明没有可读取的结果，不能证明执行从未发生。

本书新增的离线日志把顺序收紧为三个步骤。首先在同一份记录中固定候选文件束、待发布记忆快照、已结束的公开阶段资源账，以及最终测试套件的摘要。其次，把某个测试的派发预留写入事务并提交。只有提交成功，调用者才拿得到测试请求。最后，可信执行器的观察到达以后，再记录结果与清理状态。这里的“预留”指允许派发的机会，不等于确认候选已经开始执行；它宁可留下一个未用机会，也不从不确定状态中自动赚回一次尝试。

实现使用 SQLite 的显式事务组织这些记录。SQLite 的原子提交仍依赖文件系统锁、写入和刷新等条件，不能把启用事务理解为对任何硬件故障的绝对保证\cite{sqlite-atomic-commit}。本书实际检查的是新建本地文件上的重开、错误提交和受控进程退出，没有做断电或损坏磁盘实验。调用前先保存记录也不使外部服务与本地数据库成为一个共同事务；它解决的是恢复后不要盲目重发，而不是证明外部动作恰好执行一次。

\paragraph{恢复先处理未知，不能挑一个方便的结果。}
重新打开日志时，候选、记忆、限额、套件或公开资源摘要不一致，就拒绝继续。仍有未决派发或未确认清理时，后续测试同样暂缓，避免前一次残留执行污染下一次环境。如果原观察可以被可信地取回，就只结算原调用；如果无法取回，只有在确认清理后才能把该项结束为缺失，不能改记失败、补造成功或重新派发来填空。已经因清理不明记为缺失的测量，后来清理完成也不会升级成成功。

以下是实际运行过的\emph{离线构造例子}，数字不是模型能力成绩。计划有\calcnum{3}个最终测试：第一个收到匹配答案；第二个仅留下派发记录，重开后阻止第三个继续；显式完成虚构清理对账后，第二个保留缺失，第三个收到不匹配答案。报告因此是计划\calcnum{3}、观察\calcnum{2}、通过\calcnum{1}、缺失\calcnum{1}。因为已有一个明确失败，“全部通过”可以判为否；但记录仍不完整，缺失数继续保留。若已观察项全部通过而仍有缺失，“全部通过”才保持未知。不能把一比二当作完整计划的成绩，也不能把未决项补成零。公开测试账与最终测试预留分别列出，不能把同一项目重复记入两种费用。

\paragraph{可恢复日志的权限边界。}
新日志只允许追加事件，并在重开时检查事件顺序与摘要；这有助于发现意外更改，却不是防止宿主所有者重写历史的密码学保证。删除日志尾部、复制旧数据库或同时更改内容与摘要，仍需要外部封存与访问控制来防范。公开资源账和最终日志目前也是两个文件：读取前复核摘要能发现变化，不能替代跨文件原子冻结。正式运行器必须停止公开阶段写入，再进入最终评价。

每次使用相同编号创建另一份日志，也不能自动被这一局部组件发现。因此，全局运行登记、实际进程清理证明、干净实例和真实资源适配仍需接通。当前新增\calcnum{27}项检查只说明这些受控恢复路径符合合同；其中一个受控控制程序确实在预留后退出，但没有候选程序或模型后端参与。最终日志继续拒绝向模型提供反馈，记录中的候选输出也不能被当作下一轮学习材料。

\subsection{提炼请求只由已结束的公开事件构造}

“更新只看公开材料”不能仅写在提示中。若程序先把完整回合对象序列化，再删几个名字叫作隐藏结果的字段，另一个嵌套字段仍可能带出最终评分。新的离线组件改用明确的公开投影：请求只组合当前来源任务的公开题面、文件、可写清单，以及按标识取回的公开请求、回复和测试观察。宿主验收套件与迁移阶段不能作为事件来源；不存在、其他任务和关闭时点之后的事件也不能混进同一提炼窗口。

一个任务的公开事件先按发生顺序登记，再关闭窗口，最后才构造提炼请求。重复标识、改变事件顺序和伪造窗口均被拒绝。保存材料包含实际请求字节、来源事件、窗口摘要和构造器版本，便于追查某条建议当时看到了什么。摘要不是内容审查：可信调用者若把隐藏答案误写进一条标为公开的字符串，类型检查仍不能理解并排除它。因此这条接口约束用于减少意外混用，不等于候选进程隔离，也不证明文本在语义上没有泄漏。

\paragraph{存储完整笔记，读取最近四条。}
读取适配器按发布顺序选择最近四条完整笔记，再按原顺序交给当前任务。它不使用查询或相关性打分，因此应称近期读取规则，不能叫作已经验证有效的语义检索器。准备记忆要求有一至四条合格笔记；每条正文上限为四百个 Unicode 码点。过长笔记在准入处被拒绝，不在读取时截掉句尾或否定条件。码点数量也不能直接充当 token 数或费用。

一题开始时固定全量存储摘要与选中笔记，同题多个候选共用这一快照。题后更新只能影响后续任务；一次拒绝同样耗尽该题的发布决定，不提供免费重试，也不插入空笔记挤走旧经验。若尚未记录更新决定，适配器拒绝开始下一题，避免恢复时把未完成步骤悄悄略过。完成在线序列后，迁移读取器只保留冻结记忆与任务游标，没有发布笔记的接口。

构造演示从四条手写夹具笔记开始，四道在线题各追加一条，最终存储八条、迁移读取四条；这些是教学配置及程序观察，不是八条真实模型经验。新增\calcnum{18}项边界检查与相关既有检查合计\calcnum{82}项通过，没有模型请求或候选执行。这一阶段只验证了所列接口路径，当时还没有把来源凭据与笔记发布绑定。下述独立组件补充记录关联；后续离线编排已连接局部预算与持久游标，真实后端仍未接入。原有整条序列演示仍使用其原来的全量读取设置，不能将新组件的行为追记到旧运行。

提炼、语义审核与发布分别记录来源和决定。方案中的外部 Codex 语义自审属于额外审核环节，需要登记审核者、所见材料与资源用量。当前离线检查用手写审核决定和手写初始记忆核对记录关联、发布顺序与接口约束。真实实验还需按冻结的准备流程构建$M_0$，保存材料、构建轨迹、筛选规则和最终字节身份。

\subsection{准入记录绑定的是哪一段真实回复}

保存了来源请求，不代表后来发布的笔记来自这次请求。程序可能拿错回复，也可能在审核之后重新改写候选，最后留下一个看似完整、实际对应不同文本的通过记录。新的离线准入组件把这条链拆开保存：先从已关闭的公开事件窗口取得提炼请求，再登记一次回复，从回复中按明确的起止位置选取候选，然后记录针对这一候选的审核决定，最后请求发布。

这里的起止位置按 Unicode 码点计数，并保存完整回复及其摘要。候选必须是该回复中实际存在的连续片段，不能把调用者另写的一段话悄悄附上同一个回复编号。若确实需要再次改写，就属于另一次生成和审核流程，不能沿用已有决定。来源事件标识还必须属于本次请求选择的集合；来源在集合中只证明它被提供过，不证明候选忠实概括了它，更不证明后续使用一定合适。

审核记录保存请求、回复与候选摘要，以及决定、理由、审核者类别与声明身份。类别可以是脚本夹具、AI编辑或人类，但字段不认证能力。当前装置检查只使用脚本夹具。接受决定仍受长度上限约束，拒绝保留原记忆并消耗当前题的发布决定；同一回复、候选或决定不得静默替换。

发布时保存之前与之后的记忆摘要，并将决定连接到已结束的来源任务。准备阶段的记录不能直接冒充在线更新；准备记忆的汇集与合并仍需单独实现。演示从一条手写准备夹具开始，增加一条脚本回复的摘录，最终存储\calcnum{2}条；该数字只是演示的状态变化。新增\calcnum{15}项关联检查与相关既有检查共\calcnum{97}项通过，没有后端调用或候选程序执行。

上述内存组件依赖可信宿主；哈希不替代内容审核、身份认证和执行隔离。单次准入日志把提取回复、用量、审核与快照纳入同一事务（附录\ref{app:repro}），后续离线编排又将它接到运行身份、跨单元预算、序列游标与记忆合并。

\paragraph{接通序列之后，仍要说明预算覆盖什么。}
\texttt{e01\_v2\_assembly.py}用构造材料演示上述连接：已提交单元在重开后被跳过，逻辑调用与请求尝试额度、游标从本地记录恢复，已合并笔记不再次加入。模拟墙钟计数每次调用装配函数时重新建立，跳过的已完成单元不会把前次模拟耗时带回。这里的共享预算只供U臂的提取步骤支取；准备、求解、测试与传输重试尚未纳入这份演示总账。M0也来自显式合成准备材料，不是真实模型生成的初始记忆。

完成单元的重开检查不等于任意中途崩溃都能恢复。准入日志保存的发布快照也不在每一步都等于累计经验记忆，现有检查另核对条目是否进入绑定存储。要形成真实实验，还需把完整资源阶段、实际后端、候选执行与恢复边界一起接通。低层接口仍可直接调用，组件存在和局部演示通过都不能自动变成实验就绪。
